% HOT-Jass web -- Beginner's guide: "Building your own tokamak from scratch".
% Build (from this folder):  pdflatex HOT-Jass_Beginner_Guide.tex  (run twice)
% Screenshots in fig/ were taken from the running app (2026-10-14); every
% number quoted for the worked example comes from the same runs.
\documentclass[11pt,a4paper]{article}
\usepackage[T1]{fontenc}
\usepackage{lmodern}
\usepackage{amsmath,amssymb}
\usepackage[margin=22mm]{geometry}
\usepackage{booktabs,tabularx,array}
\usepackage{graphicx}
\usepackage[table]{xcolor}
\usepackage{enumitem}
\usepackage[most]{tcolorbox}
\usepackage{fancyhdr}
\usepackage{titlesec}
\usepackage[colorlinks=true,linkcolor=magenta!70!black,urlcolor=teal]{hyperref}

\definecolor{hjmagenta}{HTML}{C400C4}
\definecolor{hjteal}{HTML}{24BDB3}
\definecolor{cellyellow}{HTML}{FFF6C8}
\definecolor{cellgray}{HTML}{ECECEC}
\definecolor{cellorange}{HTML}{FFDCB8}
\definecolor{cellred}{HTML}{FFD0D0}

\titleformat{\section}{\Large\bfseries\color{hjmagenta!80!black}}{\thesection}{0.6em}{}
\titleformat{\subsection}{\large\bfseries}{\thesubsection}{0.6em}{}
\setlength{\parindent}{0pt}
\setlength{\parskip}{0.55em}
\setlist{itemsep=0.2em,topsep=0.3em}

\pagestyle{fancy}
\fancyhf{}
\fancyhead[L]{\small HOT-Jass web -- Beginner's guide}
\fancyhead[R]{\small\thepage}
\renewcommand{\headrulewidth}{0.4pt}

\newtcolorbox{keyidea}[1][]{colback=hjteal!8,colframe=hjteal!70!black,fonttitle=\bfseries,
  title=#1,arc=2mm,boxrule=0.6pt,left=2mm,right=2mm,top=1mm,bottom=1mm}
\newtcolorbox{tip}[1][Tip]{colback=hjmagenta!5,colframe=hjmagenta!60!black,fonttitle=\bfseries,
  title=#1,arc=2mm,boxrule=0.6pt,left=2mm,right=2mm,top=1mm,bottom=1mm}
\newcommand{\cell}[2]{\colorbox{#1}{\strut\makebox[2.6em]{\small #2}}}
\newcommand{\ui}[1]{\textsf{\textbf{#1}}}
\newcommand{\shot}[2]{\begin{figure}[htbp]\centering
  \includegraphics[width=\textwidth]{fig/#1}\caption{#2}\label{fig:#1}\end{figure}}

\begin{document}

\begin{titlepage}
\centering
\vspace*{30mm}
{\Huge\bfseries\color{hjmagenta!80!black} Building your own tokamak\\[2mm] from scratch\par}
\vspace{8mm}
{\Large A beginner's guide to HOT-Jass web\par}
\vspace{18mm}
\begin{minipage}{0.82\textwidth}
\large
Seven numbers describe a tokamak well enough to estimate how hot its plasma
gets, how much current it can drive and whether it can make fusion power.
This guide starts from those seven numbers, builds a machine that does not
exist yet -- \textbf{MY-ST}, a compact spherical tokamak -- from an empty
screen, and shows how HOT-Jass tells you when an input or a result needs a
second look.
\end{minipage}
\vfill
{\large October 2026\par}
\end{titlepage}

\tableofcontents
\newpage

%=====================================================================
\section{Seven numbers make a tokamak}
%=====================================================================

A tokamak looks complicated, but for a first estimate of its performance you
need only \textbf{seven numbers}: three for its shape, two for its magnetic
field, one for the plasma it holds, and one for how well it keeps heat in.

\begin{center}
\renewcommand{\arraystretch}{1.25}
\begin{tabularx}{\textwidth}{c l X c}
\toprule
 & \textbf{Parameter} & \textbf{What it sets} & \textbf{MY-ST} \\
\midrule
$R_0$ & major radius & size; volume; how far the beam travels & 0.9 m \\
$a$ & minor radius & size; aspect ratio $A = R_0/a$ & 0.55 m \\
$\kappa$ & elongation & volume and current capacity (a taller plasma holds more) & 2.0 \\
$B_0$ & toroidal field & confinement; pressure limit ($\beta$) & 1.0 T \\
$I_p$ & plasma current & confinement; stability ($q$); density limit & 1.2 MA \\
$n_e$ & electron density & how many particles share the heat; beam stopping & $6\times10^{19}\,\mathrm{m^{-3}}$ (centre)\\
$\tau_E$ & energy confinement time & how long the heat stays in & 40 ms \\
\bottomrule
\end{tabularx}
\end{center}

The eighth ingredient, the \textbf{heating power} $P$, is not a property of the
machine: it is what you \emph{do} to it. In HOT-Jass that is mostly neutral-beam
injection (NBI), with optional ECRH and ICRH.

\subsection{What follows from the seven numbers}

Everything below is a back-of-the-envelope formula, evaluated for MY-ST. The
code solves the same physics more carefully; where it matters, HOT-Jass's own
result for the same machine (Section~\ref{sec:example}) is quoted alongside.

\paragraph{Plasma volume.}
\[
  V \simeq 2\pi^2 R_0\, a^2 \kappa = 2\pi^2 \times 0.9 \times 0.55^2 \times 2.0
  = 10.75~\mathrm{m^3}.
\]

\paragraph{Safety factor (stability).}
The field lines wind around the torus; the edge safety factor $q$ counts how
many times they go round the long way for one turn the short way:
\[
  q_{\mathrm{cyl}} \simeq \frac{5\,a^2 B_0}{R_0\, I_p[\mathrm{MA}]}\,\frac{1+\kappa^2}{2}
  = \frac{5 \times 0.55^2 \times 1.0}{0.9 \times 1.2}\times\frac{1+4}{2} = 3.50 .
\]
Below $q\approx 2$ the plasma is kink-unstable and disrupts; real machines keep
$q_{95}\gtrsim 3$. A spherical tokamak gets a bonus from its tight aspect
ratio: HOT-Jass's shaped formula gives $q_{95}=6.9$ for MY-ST.

\paragraph{Density limit.}
The Greenwald density caps the line-averaged density:
\[
  n_{\mathrm{GW}}\,[10^{20}\,\mathrm{m^{-3}}] = \frac{I_p[\mathrm{MA}]}{\pi a^2}
  = \frac{1.2}{\pi\times0.55^2} = 1.26 .
\]
MY-ST's line average is $3.2\times10^{19}\,\mathrm{m^{-3}}$, a comfortable
$n/n_{\mathrm{GW}}=0.25$.

\paragraph{Temperature: the energy balance.}
In steady state the heating replaces what leaks out. With
$n_i\simeq n_e = n$ and $T_i \simeq T_e = T$, the stored energy is
$W = 3\langle n\rangle\langle T\rangle V$, so
\[
  W = P\,\tau_E \quad\Longrightarrow\quad
  \langle T\rangle \simeq \frac{P\,\tau_E}{3\,\langle n\rangle V}.
\]
A central density $n_{e0}=6\times10^{19}$ with the default profile
$n(\rho)=n_{e0}(1-\rho^2)^{2p_n}$, $p_n=1$, averages to
$\langle n\rangle = n_{e0}/(1+2p_n) = 2\times10^{19}\,\mathrm{m^{-3}}$. With
$P=3$~MW and $\tau_E=40$~ms:
\[
  \langle T\rangle \simeq \frac{3\times10^{6}\times0.04}
  {3\times2\times10^{19}\times10.75\times1.602\times10^{-16}\,\mathrm{J/keV}}
  \approx 1.2~\mathrm{keV}.
\]
HOT-Jass gives $\langle T_e\rangle = 1.29$ and $\langle T_i\rangle = 1.13$~keV.
The centre is about three times hotter than the average: $T_{e0}=3.9$~keV.

\paragraph{Pressure limit.}
The plasma pressure is measured against the magnetic pressure,
$\beta_t = 2\mu_0\langle p\rangle/B_0^2$, and normalised to the current:
\[
  \beta_N = \beta_t[\%]\,\frac{a\,B_0}{I_p[\mathrm{MA}]} .
\]
Ideal-MHD stability needs $\beta_N\lesssim 3.5$ in a conventional tokamak and
somewhat more ($\sim5$--6) in a spherical one. With its fast beam ions MY-ST
reaches $\beta_t=5.1\%$, $\beta_N = 5.1\times0.55/1.2 = 2.3$.

\paragraph{Fusion power and gain.}
For a deuterium--tritium plasma,
\[
  P_{\mathrm{fus}} = n_D\, n_T\, \langle\sigma v\rangle(T)\; E_{\mathrm{fus}}\, V,
  \qquad E_{\mathrm{fus}} = 17.6~\mathrm{MeV},
  \qquad Q = P_{\mathrm{fus}}/P_{\mathrm{heat}} .
\]
$\langle\sigma v\rangle$ rises steeply up to $T\sim 20$~keV, so the temperature
-- that is, $P\tau_E/n$ -- decides almost everything. In a beam-heated plasma
the fast beam ions also fuse with the thermal ions (\emph{beam-target}
fusion), which at a few keV is far more than the thermal plasma makes on its own.

\begin{keyidea}[Where does $\tau_E$ come from?]
$\tau_E$ is the one number you cannot read off a drawing of the machine: it
depends on turbulence in the plasma. You can either \emph{type it in} (the
simplest start, used below) or let HOT-Jass compute it from a
\emph{scaling law} fitted to many experiments, e.g.\ for spherical tokamaks
(Kaye 2006, NSTX H-mode)
\[
  \tau_E = 4.69\times10^{-9}\; I_p[\mathrm{A}]^{0.57}\, B_0^{1.08}\,
  n_e[\mathrm{m^{-3}}]^{0.44}\, P[\mathrm{W}]^{-0.73}\ \mathrm{s},
\]
or for conventional tokamaks the ITER scaling IPB98(y,2)
\[
  \tau_E = 0.0562\; I_p^{0.93} B_0^{0.15} n_{19}^{0.41} P^{-0.69}
  R_0^{1.97} \kappa^{0.78} \varepsilon^{0.58} M^{0.19}\ \mathrm{s}
  \quad (I_p\ \mathrm{MA},\ P\ \mathrm{MW},\ \varepsilon = a/R_0).
\]
HOT-Jass keeps separate electron and ion confinement times, so its rail has
two cells, \ui{tauE,e} and \ui{tauE,i}. In this guide we type the
\textbf{same} $\tau_E$ into both.
\end{keyidea}

%=====================================================================
\section{The empty rail: NEW MACHINE}
%=====================================================================

HOT-Jass opens in the \ui{NEW MACHINE} state (the magenta button above the
presets, Fig.~\ref{fig:empty.png}): every parameter in the \ui{INPUT DATA}
rail on the left is empty and every model is set to its simplest option.
Clicking \ui{NEW MACHINE} at any time brings you back here; the eight buttons
below it load preset devices instead (seven existing machines and the PONTE
design study).

\shot{empty.png}{HOT-Jass at start-up: NEW MACHINE selected, the rail empty.
Yellow cells must be filled in; gray italic text is the default an empty cell
will use.}

The colour of each value cell tells you what HOT-Jass needs:

\begin{center}
\renewcommand{\arraystretch}{1.35}
\begin{tabularx}{\textwidth}{l X}
\toprule
\textbf{Cell} & \textbf{Meaning} \\
\midrule
\cell{cellyellow}{} & \textbf{Required} and still empty. Start is blocked until you fill it in. \\
\cell{white}{\color{gray}\itshape 1 (def.)} & \textbf{Optional}: if you leave it empty, the calculation uses the gray value shown. \\
\cell{cellgray}{\color{gray}\itshape not used} & \textbf{Not used} with the current settings (a beam that is switched off, a model you have not selected). \\
\cell{cellred}{abc} & \textbf{Invalid}: not a number, out of range, or fractions that do not add up. Blocks Start. \\
\cell{cellorange}{1.5} & \textbf{Worth checking}: allowed, but probably not what you meant (Section~\ref{sec:flags}). Never blocks Start. \\
\fcolorbox{red}{white}{\strut\makebox[2.6em]{\small 0.9}} & Changed by hand since the last preset or file was loaded. \\
\bottomrule
\end{tabularx}
\end{center}

Only a handful of cells are always required: $R_0$, $a$, $B_0$, $I_p$ and the
central density $n_{e0}$ -- plus at least one heating power above zero.
A few more become required when you choose something that needs them: the two
$\tau_E$ cells while \ui{Confinement} is \ui{Fixed tauE (input)}, and the beam
energy and tangency radius of a beam whose power is above zero. An empty power
counts as zero, that is, \emph{switched off}.

\begin{tip}
Values change colour as soon as you change something they depend on: type a
beam power and that beam's cells turn from gray \emph{not used} to yellow
(required) or gray-italic defaults.
\end{tip}

%=====================================================================
\section{A worked example: MY-ST}\label{sec:example}
%=====================================================================

MY-ST is not one of the presets: a compact spherical tokamak ($A=1.64$) of
about MAST size, heated by a single 3~MW deuterium beam.

\subsection{Step 1 -- type the seven numbers}

Open the \ui{PLASMA} section and type:

\begin{center}
\begin{tabular}{lcccccccc}
\toprule
cell & \ui{R0} & \ui{a} & \ui{kappa} & \ui{B0} & \ui{Ip} & \ui{Central n\_e} & \ui{tauE,e} & \ui{tauE,i} \\
value & 0.9 & 0.55 & 2.0 & 1.0 & 1.2 & 6e19 & 0.04 & 0.04 \\
\bottomrule
\end{tabular}
\end{center}

Click a cell, type, press Enter or click the next cell. Press \ui{Start}.
HOT-Jass refuses, and the \ui{Input check} dialog says why
(Fig.~\ref{fig:seven.png}): there is \textbf{no heating}. It also has a first
piece of advice -- more on that in Step~2.

\shot{seven.png}{Step 1: seven numbers, but no heating. The Input check lists
what blocks Start (top) and what is merely worth checking (bottom).}

\subsection{Step 2 -- add a beam (and aim it)}

Open \ui{NBI-1} and type \ui{P\_NB} $=3$~MW. The beam is now on, and two of
its cells turn yellow: \ui{E\_b}, its energy -- type 60~keV -- and
\ui{tangent R}, the radius at which the beam line passes closest to the torus
axis. The others show their defaults (a deuterium beam, all at full energy). Let us make a
typical beginner's mistake and type $R_{\mathrm{tan}}=1.5$~m.
The cell turns red (Fig.~\ref{fig:miss.png}): the plasma ends at
$R_0+a=1.45$~m, so this beam would \emph{miss it}, and HOT-Jass will not start
until you correct it. The Geometry tab shows the same: the beam line passes
outside the plasma.

\shot{miss.png}{Step 2: a beam aimed at $R_{\mathrm{tan}}=1.5$~m, outside the
plasma. The cell is red -- Start is blocked -- and the Geometry view shows the miss.}

Correct it to $R_{\mathrm{tan}} = 0.9$~m (aimed at the magnetic axis, a common
choice). Then open \ui{MODELS} and change two of the simplest-option models:

\begin{itemize}
\item \ui{Orbit model} $\to$ \ui{ST orbits -- P\_phi orbit boundary}.
  The Input check already warned (Fig.~\ref{fig:seven.png}): with $R_0/a<2$ the
  large-aspect model badly underestimates how many fast ions are lost on their
  first orbit.
\item \ui{Shine-through} $\to$ \ui{Suzuki}. The \ui{Manual} model, with its
  default fraction 0, does not trace the beam at all: it counts the whole beam
  as absorbed, wherever it goes. Suzuki calculates how much of the beam passes
  straight through.
\end{itemize}

\subsection{Step 3 -- Start, and read the results}

Press \ui{Start}. After a few seconds the result tabs fill in.

\shot{base_plasma.png}{Step 3, Plasma tab: density and temperature profiles of
MY-ST. Dashed lines are the volume averages.}

\textbf{Plasma} (Fig.~\ref{fig:base_plasma.png}): $T_{e0}=3.87$~keV,
$T_{i0}=3.40$~keV in the centre; the averages $\langle T_e\rangle=1.29$,
$\langle T_i\rangle=1.13$~keV agree with the 1.2~keV of the simple estimate in
Section~1. The orange dotted curve is the density of fast beam ions.

\textbf{Power}: of the 3~MW injected, only 1.2\% shines through and 0.2\% is
lost on first orbits; 1.67~MW heats the electrons and 1.29~MW the ions.

\shot{base_current.png}{Step 3, Current tab: the beam drives 0.49~MA and the
bootstrap effect 0.35~MA of the 1.2~MA plasma current.}

\textbf{Current} (Fig.~\ref{fig:base_current.png}): the beam drives
$I_{NB}=0.485$~MA and the pressure gradient another $I_{BS}=0.345$~MA
(bootstrap current), so 69\% of the plasma current needs no transformer.

\textbf{A first flag.} The Beam tab carries an orange note: \emph{fast ions are
12\% of all ions}. The model treats the beam ions as a small population slowing
down in the plasma; above $\sim10\%$ that assumption starts to weaken. At 12\%
it is a reminder, not an alarm.

\subsection{Step 4 -- an experiment: halve the density}

Now change \ui{Central n\_e} to $3\times10^{19}$ and press \ui{Start} again.

\shot{lown_beam.png}{Step 4, Beam tab at half the density: 11.7\% of the beam
shines through, and fast ions are now half of all ions.}

Everything moves at once (Fig.~\ref{fig:lown_beam.png}):
\begin{itemize}
\item Fewer particles stop the beam: shine-through rises from 1.2\% to
  11.7\%.
\item The same beam power now heats far fewer thermal ions: $T_{i0}$ jumps
  from 3.4 to 13.1~keV.
\item Fast ions become 49\% of all ions -- the flag now says the model is
  being stretched, and the result is a rough upper estimate.
\item The beam drives 3.7~MA in a 1.2~MA plasma ($f_{NI}=3.2$): the plasma is
  \emph{overdriven}, and the assumed current profile is no longer
  self-consistent.
\end{itemize}
The lesson: NBI needs enough density to stop the beam and to share its energy.
Put $n_{e0}$ back to $6\times10^{19}$.

%=====================================================================
\section{Making it realistic}
%=====================================================================

The run in Step~3 is a valid first estimate. To make it closer to a real
plasma, replace the simplest choices one at a time, and watch what each one
does.

\paragraph{Confinement from a scaling law.}
In \ui{MODELS}, set \ui{Confinement} to \ui{Kaye NSTX H-mode} (the spherical
tokamak scaling of Section~1); the two $\tau_E$ cells turn gray, \emph{not used}.
Tick \ui{electron-ion equipartition}, the heat exchange between electrons and
ions. The scaling gives MY-ST $\tau_E = 97$~ms rather than our guessed 40~ms.

\paragraph{Shape and profiles.}
Give the plasma some triangularity, \ui{delta} $=0.4$, and a flatter density
profile, \ui{Density peaking} $=0.5$. Keep \ui{profile-corrected 0-D} ticked
(the default for NEW MACHINE): without it the power balance assumes flat
profiles and the peaking changes little.

\shot{real_plasma.png}{MY-ST with Kaye NSTX H-mode confinement,
equipartition, $\delta=0.4$ and density peaking 0.5.}

Result (Fig.~\ref{fig:real_plasma.png}): $T_{e0}=5.5$~keV, $T_{i0}=5.9$~keV,
$\beta_N=3.8$, $f_{NI}=0.84$, fast ions 7\% -- and no flags.

\paragraph{Fuel.}
So far MY-ST burns deuterium only (the default fuel mix is pure D) and makes a
mere 2~kW of D-D fusion. Set the fuel to \ui{D fraction} $=0.5$,
\ui{T fraction} $=0.5$:

\shot{dt_fusion.png}{MY-ST with a 50:50 D-T plasma: 0.62~MW of fusion power,
92\% of it from beam ions fusing with the thermal tritium.}

$P_{\mathrm{fus}} = 0.62$~MW, $Q=0.21$ (Fig.~\ref{fig:dt_fusion.png}). 92\%
of it is \emph{beam-target} fusion: at $T\approx 6$~keV the thermal plasma is
too cold to fuse much on its own, but the 60~keV deuterons are not. This is
why beam-driven spherical tokamaks are attractive as compact neutron sources.

\paragraph{Other cells worth filling in.}
\begin{itemize}
\item \ui{Z\_eff} (default 1, a perfectly clean plasma): real plasmas have
  $Z_{\mathrm{eff}}\approx1.5$--2.
\item \ui{E, E/2, E/3 power}: a positive-ion beam source also delivers atoms at
  half and a third of the full energy, e.g.\ 80/15/5\,\%.
\item \ui{width W}, \ui{height H}: the beam's cross-section (default: a single ray).
\item \ui{Centre-post R}: the radius of the central column, if it is not $R_0-a$.
\end{itemize}

\begin{tip}[Save your machine]
\ui{Config $\to$ Save} writes all inputs, including the empty cells, to a JSON
file; \ui{Read} loads it back. The file names the device ``NEW MACHINE'';
change its \texttt{device\_name} field to, say, ``MY-ST'' and that name is shown
when you load it.
\end{tip}

%=====================================================================
\section{The ``Worth checking'' flags}\label{sec:flags}
%=====================================================================

HOT-Jass checks your inputs before every Start, and the results after every
run, for the most common beginner traps. Flags never block a run; they colour
the cell orange, appear as a notification, and are listed in the
\ui{Input check} dialog or at the top of the result tab concerned.

\subsection{Before Start}
\begin{center}
\renewcommand{\arraystretch}{1.3}
\begin{tabularx}{\textwidth}{l l X}
\toprule
\textbf{Flag} & \textbf{When} & \textbf{What to do} \\
\midrule
Kink limit & $q_{95}<2$ & Lower $I_p$, or raise $B_0$, $a$ or $\kappa$. \\
Greenwald & $\langle n_e\rangle_{\mathrm{line}} > n_{\mathrm{GW}}$ & Lower $n_{e0}$ or raise $I_p$. \\
Beam hits the post & $R_{\mathrm{tan}} <$ centre-post radius & Aim it further out. \\
Orbit model & $R_0/a<2$ with the large-aspect model & Choose an ST orbit model. \\
Peaking has little effect & peaking set, profile correction off & Tick \ui{profile-corrected 0-D}. \\
$q_0$ has no effect & $q_0 \ge q_{\mathrm{cyl}}$ & Use $q_0 < q_{\mathrm{cyl}}$ (Section~\ref{sec:surprises}). \\
$q_0$ below 1 & $q_0 < 1$ & Sawteeth keep $q_0\approx1$ in a real plasma; use $q_0=1$. \\
\bottomrule
\end{tabularx}
\end{center}
Two impossible inputs are not flags but errors (red), and Start is blocked
until you correct them: $a \ge R_0$, and a beam aimed outside the plasma,
$R_{\mathrm{tan}} > R_0+a$.

\subsection{After the run}
\begin{center}
\renewcommand{\arraystretch}{1.3}
\begin{tabularx}{\textwidth}{l l l X}
\toprule
\textbf{Flag} & \textbf{When} & \textbf{Tab} & \textbf{Meaning} \\
\midrule
Shine-through & $>25\%$ of $P_{NB}$ & Power & Beam too energetic for the density; the wall takes the power. \\
First-orbit loss & $>30\%$ of $P_{NB}$ & Power & $I_p$ too low to hold the fast ions. \\
$\beta_N$ & $>3.5$ ($>5.5$ if $R_0/a<2$) & Plasma & Ideal-MHD unstable. \\
Fast ions & $>10\%$ of all ions & Beam & The slowing-down model is stretched. \\
\bottomrule
\end{tabularx}
\end{center}

\subsection{PONTE: the demanding profile}
The \ui{PONTE} preset is a design study that pushes several limits at once.
Its users get \emph{more} flags, marked \textbf{[PONTE strict]}: tighter limits
($q_{95}<3$, $n/n_{\mathrm{GW}}>0.8$, shine-through $>10\%$, first-orbit loss
$>15\%$, fast ions $>5\%$, $\beta_N > 3.0$ or $4.5$) and extra checks: the choice
of models (fixed $\tau_E$, equipartition off, an orbit model other than
P$_\phi$, D and T beams without beam-beam fusion), a strongly hot-ion plasma
($T_{i0}/T_{e0}>2$), $T_{i0}>30$~keV, thermal ions below 70\%, fast-ion
energy above the thermal one, an overdriven current ($f_{NI}>1$), and more
than 0.5~MW of shine-through. Seeing several warnings on PONTE is expected:
they show where that design is bold.

%=====================================================================
\section{Common surprises}\label{sec:surprises}
%=====================================================================

\begin{description}[style=nextline,leftmargin=0pt]
\item[``I changed $q_0$ and nothing happened.'']
$q_0$ is the on-axis safety factor of an \emph{assumed} current profile
$j\propto(1-\rho^2)^\nu$ with $\nu = q_{\mathrm{cyl}}/q_0-1$. It only affects
the bootstrap current (Current tab) and, for the P$_\phi$ orbit model or tilted
beams, the fast-ion pitch. Once $q_0\ge q_{\mathrm{cyl}}$, $\nu=0$ -- the
profile is flat -- and larger $q_0$ changes nothing. An empty $\kappa$ (i.e.\ 1)
makes $q_{\mathrm{cyl}}$ small, so this happens easily.

\item[``Density peaking does nothing.'']
With \ui{profile-corrected 0-D} off, the power balance uses flat profiles; the
peaking then only shapes the deposition and current profiles.

\item[``Shine-through is zero, whatever I do.'']
That is the \ui{Manual} shine-through model: it uses the fixed fraction you
type (default 0) and never traces the beam. Choose Suzuki, Janev or Riviere.

\item[``Why are some empty cells gray, not yellow?'']
They have a default (shown in gray italics) or are not used by the current
settings. Every default the calculation actually used is listed in the Input
check dialog and in a notification at Start.

\item[``The Current tab says $f_{NI}>1$.'']
The beam and bootstrap currents exceed $I_p$. The calculation does not feed the
current back into the plasma, so this is a sign to lower the beam power or
raise $I_p$, not a solution.
\end{description}

%=====================================================================
\section{Where to go next}
%=====================================================================
\begin{itemize}
\item Load the presets -- MAST-U, NSTX-U and ST40 are spherical tokamaks like
  MY-ST; compare their numbers with yours.
\item \ui{Config $\to$ Assumptions} lists every physics assumption and formula
  the calculation uses.
\item \ui{Refs} links the papers behind the models and the validation data.
\item \ui{Results $\to$ Save} exports a PDF report, U-files and an NBI\_OUT
  file for transport codes such as ASTRA.
\end{itemize}

\end{document}
